\scanpage{13}
\mspage{10}
(3.1), $c_4=0$ si et seulement si $a_2^2=0$. Supposons qu'on ait même $a_2=0$
\[
y^2=x^3+a_4x+a_6\qquad \pi=a_4^{-1}\,dy=dx/2y\tag*{(5.1)}
\]
\[
c_4=c_6=0\qquad \Delta=-a_4^3\qquad .
\]

Les autres systèmes $(\pi,x,y)$ donnant une équation de cette forme s'obtiennent en faisant $s=t=0$ dans (1.2) et
\[
\begin{aligned}
u^4a'_4&=a_4\\
u^6a'_6&=a_6+ra_4+r^3\qquad .
\end{aligned}\tag*{(5.2)}
\]

b) $p=2$

En caractéristique 2 , on a $c_4=a_1^4$. Supposons qu'on ait $a_1=0$ ; $\pi$ étant donnée, on peut alors choisir $x$ et $y$ tels que (1.3) soit
\[
y^2+a_3y=x^3+a_4x+a_6\qquad\pi=\frac{dx}{a_3}=dy/\raisebox{-.45ex}{$x^2+a_4$}\unsure{:|!}\tag*{(5.3)}
\]
\[
b_2=b_4=0\qquad b_6=a_3^2\qquad b_8=a_4^2\qquad\Delta=a_3^4\qquad c_4=c_6=0\, .
\]
Les autres systèmes $(\pi,x,y)$ donnant une équation de cette forme s'obtiennent en faisant $r=s^2$ dans (1.2) et
\[
\begin{aligned}
u^3\unsure{a}'_3&=a_3\\
u^4a'_4&=a_4+sa_3+s^4\\
u^6a'_6&=a_6\unsure{-|\cdot}+s^2a_4+ta_3+s^6+t^2\qquad .
\end{aligned}\tag*{(5.4)}
\]
% archive sheet number: 12
